%!TEX program = lualatex
\documentclass[ngerman,headheight=70pt]{scrartcl}
\usepackage[ngerman]{babel}
\usepackage{blindtext}
\usepackage[utf8]{luainputenc}
\usepackage{scrlayer-scrpage}
\usepackage[T1]{fontenc}
\usepackage{fontspec}
\usepackage{multicol}
\usepackage{enumitem}
\usepackage[german=quotes]{csquotes}
\usepackage[bottom=2cm,footskip=8mm]{geometry}
\setmainfont{Source Sans Pro}
\MakeOuterQuote{"}

\setlength{\parindent}{0pt}
\setlength{\parskip}{6pt}

\pagestyle{plain.scrheadings}
\thispagestyle{scrheadings}
\lohead{\vspace*{-1.5cm}\includegraphics[scale=0.4]{../images/stupa_logo_rot.pdf}}
\rohead{\\Präsidium des\\Studierendenparlaments}

\lofoot[Universität Hamburg · Präsidium des Studierendenparlaments\\\hrule Von-Melle-Park 5 · D-20146 Hamburg]{Universität Hamburg · Präsidium des Studierendenparlaments\\\hrule Von-Melle-Park 5 · D-20146 Hamburg}
\cofoot[]{}
\cohead[\thepage]{}

\renewcommand{\thesubsection}{TOP \arabic{subsection}}
\renewcommand{\thesubsubsection}{\arabic{subsubsection}.}
\setcounter{subsection}{-1}

\begin{document}
    UHH · StuPa-Präsidium · Von-Melle-Park 5 · D-20146 Hamburg

    %%% PAGE ONE (EDIT LATER) %%%%
    % Hier den Titel des Protokolls eintragen
    \section*{Protokoll der 10. Sitzung des Studierendenparlaments vom 24. November 2016}

    % Protokollanten eintragen
    \textbf{Protokoll: Gunhild Berdal}\\
    \textbf{Ort: VMP5 0079}\\
    \textbf{Beginn: 20.36 Uhr}\\
    \textbf{Ende: 21.30 Uhr}

    % Wohl die komplizierteste Stelle im Protokoll
    % Als Faustregel sollten nicht mehr als 3 Namen auf einer Zeile stehen
    % Einfach an dieser Vorlage orientieren
    % Im Idealfall ist hier nicht viel zu ändern, da kaum jemand fehlt.
    % Rücktritte müssen natürlich pro Protokoll angepasst werden
    \vspace{0.5cm}
    \begin{tabular}{ll}
        Anwesend: & \\
            RCDS (5 Sitze): &  \\
             CampusGrün (14 Sitze): & Laura Franzen, Geoffrey Youett, Philipp Droll, \\
                                   & Christian Lagod, Jim Martens, Svenja Horn, \\
                                   & Martin Sievert, Armin Günther \\
             BGZM-Liste (2 Sitze) : & \\
             WiWi (2 Sitze): & \\
             Unicorns (5 Sitze): &  Johannes Peplow, Andreas Hartkamp \\
             Liste LINKS (3 Sitze): & Gunhild Berdal, Till Petersen, Sinah Mielich \\
             HWP (2 Sitze): &  \\
             MIN (4 Sitze): &  \\
             SDS* (3 Sitze): & Mena Winkler, Jacob Petersein, Artur Brückmann \\
             Bart-LISTE (2 Sitze): & Timo Zeimet \\
             LHG (1 Sitz): & Tobias Heisig \\
             harte zeiten (1 Sitz): & Tobias Berking \\
             Jura (1 Sitz): & \\
             AL (2 Sitze): &  \\
            & \\
        Entschuldigt: & Ramon Weilinger, Antonia Niecke, Ramin Shakiba, \\
                      & Jennifer Maack, Benjamin Welling (RCDS), Elena Rysikova, \\
                      & Yasemin Günther, Tahnee Herzig, Mario Moldenhauer, \\
                      & Dean Haltmeyer (CG), Jakob Pape, Karen Martirosian (BGZM), \\
                      & Claas-Friso Hente, Charleen Lorenz (WiWi), Katharina Kucza, \\
                      & Annkathrin Löffler, Marielle Hermstrüwer (Unicorns), Ailina Salten,\\
                      & Lotte Rullkötter, Jan Detampel (MIN), Dominic Laumer (BART), \\
                      & Karima Schulze, Henri Weber (AL)\\
                      &\\
        Unentschuldigt abwesend: & Melf Johannsen (CG), Hakan Tasdemir, Gökce Celik (HWP),\\
                                & Nabil Jabareen (MIN), Maurice Martens (Jura) \\
                                &\\
        Rücktritte: & \\
        Mandatsverluste: & \\
    \end{tabular}
    \newpage
    %%% Vorgeschlagene Tagesordnung %%%
    \underline{Vorgeschlagene Tagesordnung}
    \begin{enumerate}[label={\textbf{Top \theenumi}},leftmargin=*]
        \item Aktueller Stand StuPa-Wahl (V1617-045) (30 Min.)
        \item Tag gegen Rechts (V1617-043, V1617-044) (45 Min.)
        \item Verschiedenes (5 Min.)
    \end{enumerate}

    \newpage

    % Hier beginnt Teil A. Bei Fortsetzungen kommt hier nur der Teil über
    % die Feststellung der endgültigen Fassung der Tagesordnung.
    % Für sonstige Sitzungen ist dies bereits das Gerüst für die Formalia.
    % Man muss nur noch die protokollierten Inhalte eintragen.

    {\Large \underline{Diskussion über Studienreform}}

    \textit{Sinah Mielich leitet die Diskussion mit einigen Thesen ein:}

    Die Studienreform sei sehr weit gediehen. Die Dogmen der Bologna-Reform -
    Wettbewerbsfähigkeit, nur der Markt regiere und die Gesellschaft sei nicht
    gestaltbar - würden zunehmend in Frage gestellt. Der Widerspruch bestehe
    zwischen Anpassung und Emanzipation.

    Ein Erfolg der Studienreform sei, diesen Gesamtzusammenhang zu begreifen und
    inhaltlich wie kulturell eine emanzipatorische Praxis zu entwickeln. Zunächst
    über die Lockerung formaler Restriktionen wie Fristen, Module, Anwesenheit
    und Prüfungen, dann verstärk inhaltlich, wie z.B. in der Überwindung von ABK
    (Allgemeinbildender Kompetenz bzw. Allgemeinen Berufsqualifizierenden
    Kompetenzen) und der Entwicklung von Projektstudium.

    Gelungen sei dies beispielsweise über die Dies Academici. Positiver Maßstab
    sei dabei die Bildung verantwortlicher Persönlichkeiten und die Lösung
    "epochal-typischer Probleme". Es gehe gesellschaftlicher und solidarischer
    zu.

    Jetzt gehe es darum, diesen Maßstab uniweit zu verallgemeinern, eine
    Möglichkeit dafür sei die Einführung einer Rahmenprüfungsordnung. Im
    Konflikt darum gehe es neu darum, ob die Universität Humankapital oder
    Citoyens herausbilden sollte, viele wollten ihre Anpassung retten, viele
    seien verunsichert. Die Frage sei, wie wir als VS gegen diese
    Verunsicherungen umso mehr die positiven Maßstäbe zur Geltung bringen können.

    \textit{In der Diskussion werden verschiedene Aspekte hervorgehoben:}

    Die Polarisierung, wie sie in der Wahl zwischen Trump und Sanders zum
    Ausdruck komme, spiegele sich auch in der Uni.

    Die Mehrheit habe kein Bock auf Bologna, es gehe um Gesamtverantwortung und
    das Studium selbst gestalten.

    Gegen die Befürchtungen von Profs und anderen vernetzen und ermutigen, aus
    dem eigenen Umkreis heraustreten und inhaltlich arbeiten.

    In manchen Bereichen exisistiere das Positive, das vermeintlich gegenüber der
    RPO verteidigt werden müsse, gar nicht. Worum geht’s in der Wissenschaft,
    sei die gefährlichste Frage.

    Ein Vorbehalt gegen die RPO von Profs sei, damit würden sie den Studierenden
    absprechen, zu studieren. Alles müsse praktisch orientiert sein, das sei in
    der Wirtschaft ist auch so.

    Auch auf dem letzten Dies sei der Konflikt gewesen: Wenn wir uns wirklich für
    Recht auf Arbeit, Recht auf Bildung usw. positionieren, dann… Es werde nicht
    offen dagegen gesprochen, sondern gedroht. Vorbehalt bis in linke Kreise
    hinein: Wenn wir nicht mitmachen, kommt es noch schlimmer.

    Instrumentell mit dem Gesellschaftlichen zu argumentieren, um eigentlich
    Freiräume durchsetzen zu wollen, bleibe defensiv, die Konkurrenz werde damit
    als ewig postuliert. Studienreform sei ein persönliches Bedürfnis.

    In einigen Bereichen seien die Diskussionskultur schlecht und die Probleme
    angeblich weit weg. Viele wollten dennoch Verbesserungen, hätten aber Angst
    vor mehr Arbeit oder Misserfolg.

    Bei der (Abschaffung der) Anwesenheitspflicht hätten viele Angst, dass keiner
    mehr hingeht, weil die Belastung so hoch ist. Viele meinten, man müsse zum
    Studieren gezwungen werden.

    Widerspruch: Es gebe gleichzeitig Befürwortung von Veränderung und Angst vor
    Veränderung, weil die Erfahrung mit Veränderungen in den letzten 20 Jahren
    überwiegend negativ sei. Die Irrationalität in der Bekämpfung der
    Studienreform sei ein Zeichen, wie weit wir sind.

    Anschließend berichten Franziska und Geoffrey von den Vorbereitungen des
    Studienreformtags am 30. November.


    {\Large \underline{Teil A}}

    \subsection{Formalia}

    \textit{Jim eröffnet um 20.36 Uhr die Sitzung.}

    \subsubsection{Geschäftsbericht Präsidium}

    Am Freitag war die Listenplatzauslosung, alle 19 Listen haben es geschafft,
    ihre Unterlagen rechtzeitig einzureichen.

    Es gab einen Einspruch gegen MIN-Liste, der vom Präsidium mehrheitlich
    abgelehnt wurde.

    Die weiteren Wahlvorbereitungen laufen auf volle Touren, es haben sich über
    fünf Leute als Wahlhelfer*innen beworben. Für die Wahlzeitung wurde eine
    Anwesenheitsstatistik angefertigt. Ein Problem bei der Urnenwahl wird die
    eingeschränkte Foyer-Nutzung durch die Brandschutzbedingungen sein, besonders
    für Versammlungen. Das Präsidium verhandelt mit der Uni-Verwaltung, die sich
    konstruktiv zeigt, vernünftige Lösungen zu finden.

    Die Beschlüsse des Studierendenparlaments zum Hamburger Haushalt wurden an
    die Bürgerschaftsfraktionen und Ausschüsse verschickt.

    Die kleine Listenvorstellung für die Briefwahlbroschüre muss bis morgen,
    Freitag, den 25. November um 18 Uhr eingereicht werden.

    Timo Zeimet berichtet, dass der FSR Biologie gewählt hat und die amtierende
    Liste gewählt wurde.

    \subsubsection{Anfragen an das Präsidium}

    Till: Hat sich am Format der Listendarstellung was geändert? Bitte achtet
    darauf, dass alle Namen auf dem Stimmzettel gleich und klein und die
    Listennamen korrekt geschrieben werden.

    Golnar: Gibt es eine Handhabe dafür, wie mit sehr ähnlichen Listen-Namen
    umzugehen ist?

    Till: Die Namen sollten die Liste bezeichnen und Unterscheidbarkeit von
    anderen Listen gewährleisten. Wie wäre es mit einem Handbuch für solche
    Erfahrungen bei den Wahlen?

    \textbf{Antworten}

    Am Format hat sich nichts geändert, die veränderte Beschreibung soll einen
    reibungslosen und schnellen Druck ermöglichen.

    Die Anregung, entwickelte Maßstäbe für den Umgang mit Einsprüchen besonders
    bezüglich ähnlicher Listennamen zu dokumentieren, wird vom Präsidium positiv
    aufgegriffen.

    \subsubsection{Geschäftsbericht AStA}

    Es gibt Berichte zum Akademischen Senat (Franziska), zum Umzug des Phil-Turms
    (Geoffrey), zu den Aktivitäten gegen G20 (Armin) und zur Konferenz
    "Lernfabriken meutern" mit Schüler*innen und Lehrer*innen (Phillip) sowie
    zu einer Veranstaltung zum zivilen Widerstand in Palästina, die nicht
    stattfinden konnte, weil ein Referent nicht ausreisen durfte (Golnar).

    \subsubsection{Anfragen an den AStA}

    Till an Golnar: Gibt es Konsequenzen daraus, dass die Veranstaltung nicht
    stattfinden konnte?

    Golnar: Es wird mit den Gruppen der Palästina-Solidarität weiter
    zusammengearbeitet, um Veranstaltung zum Thema zu ermöglichen.

    \subsubsection{Dringlichkeitsanträge des AStA}

    Keine.

    \subsubsection{Aktuelle Stunde}

    Es liegt kein Antrag vor.

    \subsubsection{Feststellung der endgültigen Fassung des Teils B der Tagesordnung}

    %%% BESCHLOSSENE TAGESORDNUNG %%%
    \underline{Endgültige Fassung der Tagesordnung}
    \begin{enumerate}[label={\textbf{Top \theenumi}},leftmargin=*]
        \item Aktueller Stand StuPa-Wahl (V1617-045) (30 Min.)
        \item Tag gegen Rechts (V1617-043, V1617-044) (45 Min.)
        \item Verschiedenes (5 Min.)
    \end{enumerate}

    Es gibt keinen Widerspruch zur vorgeschlagenen Tagesordnung.

    \subsubsection{Feststellung der Beschlussfähigkeit}

    \textit{Mit 19 Parlamentarier*innen ist das Parlament nicht beschlussfähig.}

    \textit{Die Sitzung wird um 21.30 abgebrochen.}

    Abschließend bittet das Präsidium bittet um Ideen und Anregungen für die
    Listenvorstellung am 15. Dezember, aufgrund der eingeschränkten
    Nutzungsmöglichkeiten der Foyers.

    Es wird sich dafür ausgesprochen, einen geeigneten Platz für das
    Zusammenkommen möglichst aller Listen und der Öffentlichkeit zu finden.

\end{document}
